% Canonical substantive source for Session 1. The live PDF is an automatic filtered export.

\LiveAndFullFrame{S1-01}{Structure of the course}{
\begin{itemize}
\item \key{Email:} \href{mailto:amedeo.andriollo@dauphine.psl.eu}{amedeo.andriollo@dauphine.psl.eu}
\item \key{Assessment:} 150-minute final in-class examination.
\item \key{Office hours:} feel free to (DO) come.
\item Slides/materials are updated regularly. Report a typo/mistake: corrections help the class.
\item \key{This session:} a first-pass sell-side target study before valuation, with AI as augmentation.
\item \key{Key references:} i) Rosenbaum \& Pearl, Chapter 1 ``Study the Target'' (Exhibit 1.3); ii) official OpenAI and Anthropic guidance; iii) Anthropic AI Fluency.
\end{itemize}
%\source{Instructor course information; source locators in \texttt{SLIDE\_SOURCE\_MAP.md}.}
}

\LiveAndFullFrame{S1-02}{Learning objectives}{
\begin{columns}[T,onlytextwidth]
\column{0.49\textwidth}
\begin{block}{During your internship...}
Before a sell-side team values a company, it turns a compact, incomplete pack into a \textbf{concise view} of the business, operating history, unresolved tensions and next questions.
\end{block}
\column{0.49\textwidth}
\begin{block}{Today's learning sequence}
\begin{enumerate}
\item Study TargetCo without a prompting method.
\item Compare outputs and (partial) feedback.
\item LLM and disciplined workflow.
\item Redo the task in a fresh chat, then inspect the full answer and workflow.
\end{enumerate}
\end{block}
\end{columns}
\smallrule
\key{Structure of our course:} target-study note $\rightarrow$ audited Excel workbook/change log $\rightarrow$ valuation analysis $\rightarrow$ reconciled Python/Excel analysis $\rightarrow$ checked deal-extraction table + derived transaction indicators $\rightarrow$ strategic buy-side briefing/workflow.
}

\LiveAndFullFrame{S1-03}{Study the target before choosing how to value it}{
\begin{block}{The finance operation}
``Study the Target'' means learning the company's story before selecting or benchmarking comparables. \\
A compact target study is neither due diligence nor valuation. It creates the working frame for later choices.
\end{block}
\begin{columns}[T,onlytextwidth]
\column{0.48\textwidth}
\key{Business profile}
\begin{itemize}
\item sector and subsector;
\item products and services;
\item customers and end markets;
\item distribution channels;
\item geography.
\end{itemize}
\column{0.48\textwidth}
\key{Financial profile (selected dimensions)}
\begin{itemize}
\item size;
\item profitability;
\item growth;
\item (returns and credit profile later).
\end{itemize}
\end{columns}
\sourcewithref{Rosenbaum \& Pearl (2013), Ch. 1, Step I ``Study the Target'' (Exhibit 1.3).}{S1-36}
}

\LiveAndFullFrame{S1-04}{Simple measures answer different economic questions}{
\begin{columns}[T,onlytextwidth]
\column{0.49\textwidth}
\begin{block}{Revenue growth and its disclosed bridge}
\[
g_t=\frac{R_t}{R_{t-1}}-1,\qquad g_t=v_t+p_t.
\]
Here $v_t$ and $p_t$ are disclosed volume and price/mix \textbf{contribution points} to revenue growth.\\
\textbf{Question answered:} how did sales change, and which disclosed contribution drove the change?
\end{block}
\column{0.49\textwidth}
\begin{block}{Reported EBITDA margin}
\[
m_t=\frac{\text{EBITDA}_t}{R_t}.
\]
\medskip
\textbf{Question answered:} how much operating profit is reported per euro of revenue?
\end{block}
\end{columns}
\smallrule
Growth vs. margin: they can move differently, as margin also depends on cost recovery, utilization, mix and other operating factors. \\ 
\key{TargetCo:} 2025A utilization exceeds 2022A, while the 2025A reported margin is lower. \\
$\hookrightarrow$ analytical tension, not causality.
\sourcewithref{Rosenbaum \& Pearl (2013), Ch. 1.}{S1-36}
}

\LiveAndFullFrame{S1-05}{Baseline design: today's first assignment}{
\begin{center}
\footnotesize
\begin{tabular}{p{25mm}p{32mm}p{32mm}p{32mm}}
\toprule
 & \textbf{ChatGPT group} & \textbf{Claude group} & \textbf{Human-only group}\\
\midrule
Inputs & \multicolumn{3}{l}{Same two-page TargetCo profile; same VP assignment; same time.}\\
Baseline & Student-composed 1 message prompt from common sketch. & Same finance assignment. & Same finance assignment w/out AI.\\
Redo & Temporary chat (max 3 messages). & Incognito chat (max 3 messages). & Possible AI access.\\
Main comparison & \multicolumn{3}{l}{Baseline output/work process versus redo output/work process}\\
\bottomrule
\end{tabular}
\end{center}
\begin{alertblock}{Interpretation discipline}
The baseline is not a ChatGPT vs. Claude comparison: students may formulate different prompts.\\
$\hookrightarrow$ It records different student/AI workflows under a common finance assignment, material, time and prompt budget.
\end{alertblock}
}

\LiveAndFullFrame{S1-06}{Baseline prompt sketch}{
\begin{block}{AI groups: one message, with student freedom}
You may edit this modest starting sketch before sending your one permitted baseline message.\\ 
The VP request remains the separate finance assignment.
\medskip
\begin{quote}\small
Using the supplied TargetCo profile, help me prepare the one-page note requested by the VP. Use the supplied material as the source. Do not value the company.
\end{quote}
\end{block}
\key{Human-only group:} complete the same finance assignment from the same profile, without AI.\par\medskip
\key{Record:} retain the exact prompt your group actually used: NO second message allowed.
%\source{INSTRUCTOR SYNTHESIS: \texttt{student/TargetCo\_baseline\_prompt\_sketch.md}.}
}

\LiveAndFullFrame{S1-07}{20-min baseline: produce a first-pass VP note}{
\begin{block}{Deliverable}
One page. Use only the supplied material. Do not value TargetCo.\\ AI groups retain the actual conversation; human-only groups retain notes.
\end{block}
\begin{columns}[T,onlytextwidth]
\column{0.48\textwidth}
\key{Read for}
\begin{itemize}
\item business configuration;
\item commercial facts and operating history;
\item what is known and what is missing.
\end{itemize}
\column{0.48\textwidth}
\key{Do not do}
\begin{itemize}
\item external research;
\item peer selection or valuation;
\item a generic diligence inventory;
\item a prompting-method exercise.
\end{itemize}
\end{columns}
\vfill
\centering\Large \textcolor{CourseNavy}{Start now. We compare outputs in 20'.}
}

\LiveAndFullFrame{S1-08}{A stronger single-shot instruction}{
\scriptsize
\begin{block}{Instructor example: still one message}
Use only the supplied TargetCo profile as the permitted source. Prepare one concise, one-page first-pass target-study note for the VP before valuation. Do not use external facts and do not value the company.
\begin{enumerate}
\item \textbf{Business model / positioning:} synthesize products, qualification, customers/end markets, direct/distributor route to market and geography. State only what the profile supports.
\item \textbf{Historical trajectory:} recompute revenue growth where possible. Treat volume and price/mix as disclosed contribution points to revenue growth. Interpret utilization and reported EBITDA margin without asserting an unproven causal bridge.
\item \textbf{Key tensions:} identify the material evidence gaps or management statements that require qualification.
\item \textbf{Three ranked questions + why each matters:} give exactly three questions. For each, name the later judgement it could change: revenue, margin, normalization, risk, peer selection or valuation.
\end{enumerate}
For each material number or management statement, use a brief profile/table locator. Keep observed data, management statements and analyst inferences distinct.
\end{block}
\sourcewithref{Spoiler of the full sequence.}{S1-50}
}

\LiveAndFullFrame{S1-09}{Four observable dimensions for the VP note}{
\footnotesize
\begin{tabularx}{\textwidth}{@{}p{35mm}X@{}}
\toprule
\textbf{Dimension} & \textbf{Score 2/2:}\\
\midrule
Business-model synthesis & Connects product specification, customer/end-market mix, direct/distributor route to market and geography/footprint into a coherent operating configuration. \\
                            & It does not invent pricing power.\\
Historical economics & Reconciles revenue growth; distinguishes volume contribution from price/mix contribution; interprets utilization and reported EBITDA-margin movement; identifies association rather than unproven causality.\\
                            & Numerical fidelity belongs here.\\
Certainty vs. Uncert. & Keeps observed data and management statements distinct; qualifies causal interpretation; imports no external facts; turns unsupported statements into questions where appropriate.\\
VP usefulness & Uses exactly three ranked questions. Each says which later judgement could change. The note is concise and decision-useful.\\
\bottomrule
\end{tabularx}
\vspace{2pt}
\textbf{Rate yourself:} \risk{0} absent, incorrect; \textbf{1} present but partial, or mechanically copied; \good{2} meets standards.
\sourcewithref{Common rubric for baseline, redo and optional homework.}{S1-38}
}

\LiveAndFullFrame{S1-10}{Partial answer: an analytical route}{
\begin{columns}[T,onlytextwidth]
\column{0.48\textwidth}
\key{Route}
\begin{enumerate}
\item Build a synthesis of the operating model.
\item Explain the historical change in drivers, rather than list annual snapshots.
\item Name tensions created by the supplied doc.
\item Ask only questions whose answers could change a later valuation judgement.
\end{enumerate}
\column{0.48\textwidth}
\key{Not shown yet}
\begin{itemize}
\item complete model answer;
\item complete ranked-question set;
\item full fact/management/inference treatment;
\item multi-turn instructor workflow;
\item final source reconciliation.
\end{itemize}
\end{columns}
\begin{block}{Professional destination at the halfway point}
Make a small evidence pack decision-useful while preserving what it does not establish.
\end{block}
\sourcewithref{What the two-page Target Co profile can support}{S1-37}
}

\LiveAndFullFrame{S1-11}{Working interpretation: revenue as bridge of contributions}{
\begin{center}
\small
\begin{tabular}{lrrrr}
\toprule
 & 2022A & 2023A & 2024A & 2025A\\
\midrule
Revenue growth & -- & 2.1\% & 1.6\% & 5.9\%\\
Volume contribution (ppt) & $-2.5$\% & $-4.0$\% & 0.5\% & 4.0\%\\
Price / mix contribution (ppt) & 10.0\% & 6.1\% & 1.1\% & 1.9\%\\
Utilization & 82\% & 78\% & 74\% & 84\%\\
Reported EBITDA margin & 14.8\% & 14.0\% & 12.8\% & 14.3\%\\
\bottomrule
\end{tabular}
\end{center}
\scriptsize
\textbf{How to read the table.} Positive price/mix contribution can offset falling volume. Recovering utilization can coexist with incomplete margin recovery. The table therefore shows co-movement, not a causal margin!\par\smallskip
\begin{block}{A defensible reading}
In 2023A, +6.1 ppt price/mix offset a $-4.0$ ppt volume contribution.\\ In 2025A, +4.0 ppt volume is the larger disclosed contributor. \\
Utilization recovered to 84\%, but the 14.3\% margin remains 50 bp below 2022A. \\
Recovery is evident; the supplied pack does not quantify the margin bridge or prove a steady-state outcome.
\end{block}
\sourcewithref{Full commentary.}{S1-47}
}

\LiveAndFullFrame{S1-12}{Partial feedback: business-model synthesis}{
\begin{columns}[T,onlytextwidth]
\column{0.54\textwidth}
\begin{block}{What the pack supports}
TargetCo is a specification-led manufacturer rather than merely a reseller. Engineering, tooling and customer qualification, combined with a 71\% direct/key-account mix, indicate a service- and qualification-intensive route to market. The distributor channel and European footprint add commercial context.
\end{block}
\begin{block}{What it does not support/establish}
It does not establish pricing power, retention, channel margin or contract durability. The 9.4\% largest-customer and 41.8\% top-ten figures are facts requiring context, not an automatic risk classification.
\end{block}
\column{0.42\textwidth}
\begin{block}{Valuation-relevant question}
\textbf{What is the end-market, channel and facility-level evidence behind the 2025A volume recovery?}\\[3pt]
Its answer could change revenue confidence, risk and later peer selection.
\end{block}
\end{columns}
\sourcewithref{What the two-page Target Co profile can support}{S1-37}
}

\LiveAndFullFrame{S1-13}{LLMs generate conditional text; they do not supply a truth test}{
\footnotesize
\begin{tabularx}{\textwidth}{@{}p{29mm}X@{}}
\toprule
\textbf{Term} & \textbf{Definition}\\
\midrule
AI & Computational systems that perform tasks associated with perception, prediction, generation or decision support.\\
Generative AI & AI that produces new content, such as text, code, images or audio.\\
LLM & A generative model trained to predict and generate sequences of language tokens.\\
Token & A unit of text processed by the model; a token is not reliably a word.\\
Context window & The finite token sequence the model can use at one time: instructions, documents, prior messages and current task instruction.\\
Inference / generation & Running the already-trained model: given current context and tokens already generated, the model computes a conditional distribution over the next token; the decoding rule selects or samples a token and the process repeats until completion.\\
\bottomrule
\end{tabularx}
\begin{alertblock}{Analytical implication}
An answer can be grammatically convincing while misstating a source, skipping a number or inventing a causal story. Language quality is not source fidelity.
\end{alertblock}
\sourcewithref{Model, Application Tool in workflow}{S1-39}
}

\LiveAndFullFrame{S1-14}{Autoregressive generation explains why context changes outputs}{
\begin{block}{A minimal technical model}
\[
p_\theta(x_1,\ldots,x_T\mid c)=\prod_{t=1}^{T}p_\theta(x_t\mid x_{<t},c).
\]
\end{block}
\begin{itemize}
\item $x_t$ is the next token; $x_{<t}$ is the preceding generated sequence.
\item $c$ is supplied context: task, source material, examples and prior conversation.
\item $\theta$ denotes learned model parameters. At generation time, the application repeatedly chooses or samples a next token from a conditional distribution.
\end{itemize}
\begin{alertblock}{Consequence}
Changing a prompt or context changes the distribution of possible outputs. It does not create a criterion for whether an output is true, financially reasonable or properly sourced.
\end{alertblock}
\sourcewithref{Model, Application Tool in workflow}{S1-39}
}

\LiveAndFullFrame{S1-15}{Context design: evidence, authority and salience}{
\scriptsize
\begin{columns}[T,onlytextwidth]
\column{0.49\textwidth}
\begin{block}{A bit of terminology:}
\textbf{Evidence:} information from the (permitted) source used to support a claim/calculation.\par\smallskip
\textbf{Source authority:} the status assigned by the analyst/task specifying which source governs a claim when sources differ.\par\smallskip
\textbf{Salience:} how strongly a supplied item effectively influences generation.\\ 
Analyst importance does not imply model salience.
\end{block}
\column{0.49\textwidth}
\begin{block}{TargetCo contrast}
The profile is \textbf{authoritative for this exercise}.\\ 
M-01 and M-02 inside it remain \textbf{management statements}, not independently established facts.\par\smallskip
More context can add evidence but also irrelevant instructions, competing sources, weak salience and verification cost. Identify the permitted source, preserve its boundary, then request extraction before synthesis.
\end{block}
\end{columns}
\sourcewithref{Long-context failures}{S1-40}
}

\LiveAndFullFrame{S1-16}{Model vs. Application vs. Tool}{
\small
\begin{center}
TargetCo PDF $\rightarrow$ application handles/parses file and manages chat/context $\rightarrow$ underlying model generates response $\rightarrow$ calculator, code or file tool may retrieve, recompute or check $\rightarrow$ analyst reviews result
\end{center}
\begin{tabularx}{\textwidth}{@{}p{27mm}X@{}}
\toprule
\textbf{Object} & \textbf{Meaning in this pathway}\\
\midrule
Model & The underlying generative model that produces conditional token predictions. A displayed version can matter, but it is not the entire workflow.\\
Application & The interface/orchestration layer that may handle files, memory, system instructions, context and model routing. Comparing applications does not isolate only their models.\\
Tool & An external capability, such as a calculator, code execution or file operation. It can retrieve or recompute something, creating a separate source and checking obligation.\\
\bottomrule
\end{tabularx}
\sourcewithref{Model, Application Tool in workflow}{S1-39}
}

\LiveAndFullFrame{S1-17}{AI fluency = human work before, during and after AI use}{
\begin{columns}[T,onlytextwidth]
\column{0.48\textwidth}
\begin{block}{Anthropic's 4D framework}
\textbf{Delegation} -- decide whether, when and what to hand to AI.\par
\textbf{Description} -- communicate goal, context and standards.\par
\textbf{Discernment} -- assess quality and usefulness.\par
\textbf{Diligence} -- take responsibility for use and final delivery.
\end{block}
\column{0.48\textwidth}
\begin{block}{Modes}
\textbf{Automation:} AI performs a routine task.\par
\textbf{Augmentation:} AI expands a person's analysis while the person directs and judges it.\par
\textbf{Agency:} AI pursues a multi-step objective with delegated action scope.
\end{block}
\end{columns}
\smallrule
\key{Course map:} Session 1 is principally \textbf{augmentation}: delegate extraction and drafting, retain judgement and verification. Session 6 moves closer to agency, with plans, approval points and stopping rules.
\sourcewithref{Session 1 and AI fluency 4Ds}{S1-52}
}

\LiveAndFullFrame{S1-18}{Success criteria and evaluation}{
\scriptsize
\begin{block}{Being precise:}
\textbf{Success criterion:} property the output must satisfy.
VS. \textbf{Evaluation:} observable test. \\
VS. \textbf{Failure example:} result showing the criterion was not met.
\end{block}
\begin{tabularx}{\textwidth}{@{}p{29mm}p{48mm}X@{}}
\toprule
\textbf{Criterion} & \textbf{Evaluation} & \textbf{Failure example}\\
\midrule
Historical economics & Check that 2023A says +6.1 ppt price/mix offsets $-4.0$ ppt volume and reconciles the displayed 2.1\% growth. & Calls 2023A volume-led growth.\\
Evidence discipline & Inspect the treatment of M-02. & States commissioning costs are non-recurring as an established fact.\\
Business-model synthesis & Check whether qualification, route to market and footprint form an operating configuration. & Calls TargetCo a reseller or asserts pricing power.\\
Prioritization / VP usefulness & Count exactly three ranked questions and inspect the named later judgement. & Gives an unranked generic diligence list.\\
\bottomrule
\end{tabularx}
\sourcewithref{Evaluation comes before prompt cleverness.}{S1-41}
}

\LiveAndFullFrame{S1-19}{Prompt ablation: P0 through P4}{
\scriptsize
\begin{tabularx}{\textwidth}{@{}p{13mm}p{46mm}X@{}}
\toprule
\textbf{Prompt} & \textbf{Increment} & \textbf{What it is intended to change}\\
\midrule
P0 & Student-composed baseline prompt from common sketch. & Establish the baseline workflow without an advanced method.\\
P1 & (Strong) Instructor's single-shot instruction, shown after baseline. & Makes the permitted source and four output sections explicit; requests exactly three ranked questions.\\
P2 & Analysis before drafting. & Makes extraction, contribution arithmetic and inference reviewable before prose hides an error.\\
P3 & Explicit criteria and checks. & Uses the four observable success criteria on this slide and checks material source treatment.\\
P4 & Extract, challenge, draft and verify across a multi-turn workflow. & Localizes errors and creates human handoff points.\\
\bottomrule
\end{tabularx}
\begin{alertblock}{Method}
If a response uses information outside the permitted source or answers a different task, improve the specification and source boundary. Missing evidence or an unsuitable tool cannot be fixed merely by longer wording.
\end{alertblock}
\sourcewithref{Prompt-chain skeleton.}{S1-45}
}

\LiveAndFullFrame{S1-20}{P1 and P2: specification before polished prose}{
\begin{columns}[T,onlytextwidth]
\column{0.49\textwidth}
\begin{block}{P1: strong single-shot instruction}
Use only the supplied TargetCo profile. Produce four sections: business model / positioning; historical trajectory; key tensions; three ranked questions + why each matters. No external facts or valuation.
\end{block}
\begin{block}{What P1 changes}
It makes source authority and the requested output structure visible. It still leaves the reasoning process hidden.
\end{block}
\column{0.49\textwidth}
\begin{block}{P2: analysis before draft}
Before drafting, extract business-model evidence; recompute historical growth; distinguish volume and price/mix contribution points; track utilization and margin; list what each point supports and does not establish. No final prose yet.
\end{block}
\begin{block}{What P2 changes}
It creates a reviewable intermediate artifact. It does not make the model's interpretation independently true.
\end{block}
\end{columns}
\sourcewithref{Prompt-chain skeleton.}{S1-45}
}

\LiveAndFullFrame{S1-21}{P3 and P4: an explicit test and a reviewable chain}{
\begin{columns}[T,onlytextwidth]
\column{0.49\textwidth}
\begin{block}{P3: attach criteria and checks}
Check the four observable success criteria: operating configuration; historical economics including contribution arithmetic; evidence/uncertainty discipline; exactly three ranked, valuation-relevant questions in a concise VP note.
\end{block}
\begin{block}{P4: make a chain}
\textbf{Extract evidence} $\rightarrow$ \textbf{challenge against criteria} $\rightarrow$ \textbf{draft VP note} $\rightarrow$ \textbf{reconcile to source}.
\end{block}
\column{0.49\textwidth}
\begin{block}{Why chaining is a design choice}
\textbf{Advantages:} intermediate outputs are reviewable; errors are easier to localize; humans can intervene.\par\smallskip
\textbf{Costs:} more interactions and latency; state/context management; new opportunities for inconsistency between steps; paid vs. unpaid versions.
\end{block}
\end{columns}
\sourcewithref{Prompt-chain skeleton.}{S1-45}
}

\LiveAndFullFrame{S1-22}{Examples, structure and roles}{
\footnotesize
\begin{tabularx}{\textwidth}{@{}p{29mm}X X@{}}
\toprule
\textbf{Technique} & \textbf{What it changes and when it helps} & \textbf{What it cannot do}\\
\midrule
Examples & When the form is ambiguous: a non-TargetCo output-shape example can demonstrate phrasing or a compact question format.  & It cannot (and should not) supply TargetCo evidence. A filled TargetCo example before baseline leaks the analytical answer.\\
Structured prompt & Clear headings separate task, sources, output and checks when a prompt has several components. & It cannot validate numbers or make a management statement true.\\
Role & ``Write for a VP from the supplied profile'' can set audience, tone and evidence discipline. & It cannot provide privileged facts, prove causality or transfer accountability.\\
\bottomrule
\end{tabularx}
\begin{block}{Concrete role comparison}
\textbf{Weak:} ``You are the world's best banker. Tell me what will happen.''\\
$\hookrightarrow$ \textbf{Useful:} ``Write for a VP from the supplied material; label uncertainty rather than invent detail.''
\end{block}
\sourcewithref{Examples of Appropriate vs. Inappropriate.}{S1-42}
}

\LiveAndFullFrame{S1-23}{Source architecture and structured prompts}{
\scriptsize
\begin{block}{Definitions}
\textbf{Task instruction} = the question or work order the model must answer after receiving documents.\quad
\textbf{Source-grounded extraction} = return the relevant fact, statement or calculation with its source location before asking the model to synthesize across documents.
\end{block}
\begin{columns}[T,onlytextwidth]
\column{0.47\textwidth}
\key{Neutral portable form}\par\smallskip
\ttfamily
Task: prepare a one-page TargetCo VP note.\\
Sources: supplied TargetCo profile only.\\
Output: four named sections.\\
Criteria: four observable success criteria.\\
Checks: recompute revenue bridge; reconcile material claims.
\column{0.50\textwidth}
\key{Anthropic XML form}\par\smallskip
\ttfamily
\textless task\textgreater prepare a one-page TargetCo VP note \textless/task\textgreater\\
\textless sources\textgreater supplied TargetCo profile only \textless/sources\textgreater\\
\textless output\textgreater four named sections \textless/output\textgreater\\
\textless criteria\textgreater four observable success criteria \textless/criteria\textgreater\\
\textless checks\textgreater recompute bridge; reconcile claims \textless/checks\textgreater
\end{columns}
\smallskip
Both express the same logical structure. 
\sourcewithref{Example of neutral vs. anthropic.}{S1-43}
}

\LiveAndFullFrame{S1-24}{Verification hierarchy}{
\small
\begin{tabularx}{\textwidth}{@{}p{31mm}p{52mm}X@{}}
\toprule
\textbf{Level} & \textbf{Question} & \textbf{TargetCo example}\\
\midrule
1. Output / completeness & Did the draft use the four sections and exactly three questions? & Count sections and questions; check one-page form.\\
2. Source fidelity & Does every material number and statement reconcile to supplied evidence? & Confirm 2025A revenue = EUR 481.0m; preserve M-02 as management's statement.\\
3. Independent check & Can a calculation, source or tool not dependent on the generated prose verify the result? & Recompute $481.0/454.0-1=5.9\%$; inspect the source table or a separate spreadsheet.\\
\bottomrule
\end{tabularx}
\begin{alertblock}{Boundary}
Asking the same model ``are you sure?'' may detect an omission. It is not an independent verification.
\end{alertblock}
\sourcewithref{Self-check vs. independent check.}{S1-46}
}

\LiveAndFullFrame{S1-25}{How AI can help make the distinction operational}{
\begin{columns}[T,onlytextwidth]
\column{0.31\textwidth}
\begin{block}{Observed data}
Supplied operating data or source content.\par\smallskip
\textbf{Example:} 2025A utilization = 84\%.
\end{block}
\column{0.34\textwidth}
\begin{block}{Management statement}
What management says; it is informative but not independently established.\par\smallskip
\textbf{Example:} M-02 calls two years of costs non-recurring.
\end{block}
\column{0.31\textwidth}
\begin{block}{Analyst inference}
A qualified interpretation drawn from facts and statements.\par\smallskip
\textbf{Example:} recovery is evident, but the pack does not prove a steady-state margin.
\end{block}
\end{columns}
\begin{block}{Operational workflow}
1. AI proposes a classification for each material statement. \\
2. AI supplies a source locator and rationale. \\
3. Analyst verifies source and decides whether inference is warranted.\\
4. Correct misclassified or unsupported claims before drafting.
\end{block}
%\sourcewithref{INSTRUCTOR SYNTHESIS; TARGETCO APPLICATION.}{S1-54}
}

\LiveAndFullFrame{S1-26}{An evidence matrix records support and limits}{
\scriptsize
\begin{tabularx}{\textwidth}{@{}p{29mm}p{33mm}p{28mm}X@{}}
\toprule
\textbf{Claim / source location} & \textbf{Observed data, management statement or calculation} & \textbf{Supports} & \textbf{Does not establish / open question}\\
\midrule
2025A growth character, Historical table & Volume contribution +4.0 ppt; price/mix contribution +1.9 ppt; revenue +5.9\% & 2025A differs from 2023A & Breadth by customer, channel or end market\\ \\
Utilization/margin tension: 2022A and 2025A & Utilization 82\% $\rightarrow$ 84\%; margin 14.8\% $\rightarrow$ 14.3\% & Utilization alone is insufficient & Mix, cost, pricing and facility bridge\\ \\
M-02 needs challenge & Management calls costs in 2024A and 2025A ``non-recurring'' & Recurrence is valuation-relevant & Whether any cost is recurring or a justified adjustment\\ \\
\bottomrule
\end{tabularx}
\vspace{2pt}
The matrix dissects: claim, source, classification/calculation, support and limit.\\
It makes uncertainty explicit rather than hiding it in polished prose.
\sourcewithref{Full(er) evidence matrix.}{S1-47}
}

\LiveAndFullFrame{S1-27}{Redo: a new chat, the same source pack}{
\begin{block}{15-minute redo}
Start a \textbf{different fresh ephemeral chat}. ChatGPT groups use Temporary chat (where available). Claude groups use Incognito chat (where available). \\
Former human-only groups receive AI access and may split between applications.
\end{block}
\begin{columns}[T,onlytextwidth]
\column{0.48\textwidth}
\key{Maximum 3 messages}
\begin{enumerate}
\item evidence / analysis;
\item challenge + draft;
\item verify + revise.
\end{enumerate}
Do not send a message to erase context.\\ 
The fresh chat separates redo from baseline.
\column{0.48\textwidth}
\key{Record}
\begin{itemize}
\item actual baseline prompt;
\item fresh-chat mode and three prompts/labels;
\item delegated versus retained judgement;
\item challenged statement; numerical/source check; remaining uncertainty;
\item short baseline-versus-redo reflection.
\end{itemize}
\end{columns}
\sourcewithref{Instructor's prompt sequence.}{S1-50}
}

\LiveAndFullFrame{S1-28}{Post-redo answer I: synthesis and historical economics}{
\begin{block}{Business model / positioning}
\footnotesize
TargetCo is a European specialty manufacturer of technically specified rigid containers, closures and dispensing assemblies. Engineering, tooling and customer qualification form part of the offer. Direct/key accounts represent 71\% of FY2025A revenue; selective distributors represent 29\%. Food \& Beverage is the largest end market (55\%), followed by Healthcare \& Pharmaceutical (24\%) and Premium Consumer (21\%); France (32\%) and DACH (24\%) are the largest disclosed geographies. The 9.4\% largest-customer and 41.8\% top-ten measures describe concentration only: the pack provides no benchmark, identities, contracts or retention evidence.
\end{block}
\begin{block}{Historical trajectory}
\footnotesize
Revenue rose from EUR 438.0m in 2022A to EUR 481.0m in 2025A. In 2023A, +6.1 ppt price/mix offset a $-4.0$ ppt volume contribution; this differs from 2025A's +5.9\% growth, where +4.0 ppt volume is the larger disclosed contributor. Utilization fell to 74\% in 2024A during Łódź commissioning and recovered to 84\% in 2025A; reported EBITDA margin recovered from 12.8\% to 14.3\%, still below 2022A's 14.8\%. Recovery is supported, but the pack does not prove its full causal bridge.
\end{block}
\sourcewithref{Full(er) evidence matrix.}{S1-47}
}

\LiveAndFullFrame{S1-29}{Post-redo answer II: tensions and ranked questions}{
\footnotesize
\begin{columns}[T,onlytextwidth]
\column{0.47\textwidth}
\key{Why these are tensions}
\begin{itemize}
\item \textbf{Utilization vs. Margin:} 2025A utilization exceeds 2022A, but margin is 50 bp lower. If utilization is a main recovery lever, an unobserved margin bridge remains.
\item \textbf{(Management-inferred) Recovery:} M-01 cites utilization, mix and efficiency, but the pack does not quantify their contributions.
\item \textbf{``Non-recurring'' across two years:} two-year duration weakens automatic one-off treatment, but does not prove recurrence.
\end{itemize}
\column{0.50\textwidth}
\key{Three ranked questions -- and why}
\begin{enumerate}
\item \textbf{What bridges margin from 2022A to 2025A and supports run-rate?}\\
It could change margin forecasts and valuation.
\item \textbf{What is the nature and duration of 2024A-2025A commissioning/quality costs?}\\
It could change normalization judgement.
\item \textbf{How broad and durable is volume recovery by customer, channel, end market and facility?} 
It could change revenue confidence, risk and peer selection.
\end{enumerate}
\end{columns}
\begin{alertblock}{Evidence limitation}
M-01 and M-02 remain management statements.\\ 
The supplied pack supports neither a forecast nor a valuation conclusion.
\end{alertblock}
\sourcewithref{Why the three questions are ranked.}{S1-48}
}

\LiveAndFullFrame{S1-30}{Instructor workflow: AI helps when its work is challengeable}{
\small
\begin{enumerate}
\item \key{Evidence / analysis:} work from the previous slide's tensions. Extract business-model evidence; recompute revenue growth; identify contribution points; track utilization and margin. Return source location, support and limit. No final prose.
\item \key{Challenge + draft:} test the analysis against four criteria, correct unsupported inference or generic questions, then draft the one-page note in four sections.
\item \key{Verify + revise:} reconcile every number and material statement to source or calculation. Revise only where the evidence supports revision; flag unresolved items.
\end{enumerate}
\begin{block}{The analyst still determines}
Source authority; whether causal language is justified; which tension matters for valuation; whether a cost qualifies for normalization; whether the final answer is acceptable.
\end{block}
\sourcewithref{Instructor's prompt sequence. Full prompt in \texttt{SESSION1\_INSTRUCTOR\_AI\_WORKFLOW.md}.}{S1-50}
}

\LiveAndFullFrame{S1-31}{Take some time to compare the actual outputs}{
\footnotesize
\begin{columns}[T,onlytextwidth]
\column{0.53\textwidth}
\key{Baseline versus redo}
\begin{itemize}
\item Which answer correctly decomposes the historical revenue story?
\item Which recognizes the utilization/margin tension?
\item Which question would change later valuation assumptions most?
\item Which output reports a management assertion as if independently verified?
\item Which output invents a causal explanation unsupported by the pack?
\end{itemize}
\column{0.43\textwidth}
\key{Optional homework: same-prompt cross-application comparison}
Use the same instructor near-final prompt sequence, same TargetCo profile and fresh Temporary/Incognito chats. Match tools, score both on the common rubric, record date/application/displayed model-version/tool-memory state, then write what differences matter to a VP.
\end{columns}
\smallskip
%\sourcewithref{INSTRUCTOR SYNTHESIS: \texttt{instructor/BASELINE\_DISCUSSION\_GUIDE.md}.}{S1-50}
}

\LiveAndFullFrame{S1-32}{Each of our sessions changes the work product: not accountability!}{
\small
\begin{tabularx}{\textwidth}{@{}p{16mm}X@{}}
\toprule
\textbf{Session} & \textbf{Professional work product}\\
\midrule
1 & Target-study note.\\
2 & Audited Excel workbook and change log.\\
3 & Valuation analysis.\\
4 & Reconciled Python/Excel analysis.\\
5 & Checked deal-extraction table plus derived transaction indicators.\\
6 & Strategic buy-side briefing/workflow.\\
\bottomrule
\end{tabularx}
\begin{block}{Continuity}
The finance artifact, technical environment and permitted AI role become more demanding. \\
$\hookrightarrow$ The analyst retains accountability for specification, evidence, judgement and verification.
\end{block}
}

\LiveAndFullFrame{S1-33}{Session 2 preparation}{
\footnotesize
\begin{block}{Next professional question}
Can we rely on TargetCo's historical Excel model? \\
The task moves from target-study narrative to financial-statement and model logic (not yet valuation).
\end{block}
\begin{columns}[T,onlytextwidth]
\column{0.49\textwidth}
\key{Refresh}
\begin{itemize}
\item reported/adjusted EBITDA; debt, cash, net debt and leverage;
\item CapEx/NWC consistency and formula tracing;
\item Rosenbaum \& Pearl, \emph{Investment Banking}, Ch. 1: ``Financial Profile'' and ``Calculation of Key Financial Statistics and Ratios.''
\end{itemize}
\column{0.49\textwidth}
\key{Book references}
\begin{itemize}
\item Workbook Ch. 1: business/financial profile Q14--21; growth/profitability/credit Q7--9; one-time/non-recurring Q3.
\item Koller, Goedhart \& Wessels (2025), Chs. 11 and 13, for optional model organisation/hygiene context.
\end{itemize}
\end{columns}
\smallskip
Use these sections to refresh the concepts before class; you are not expected to reproduce or solve the textbook exercises.
\sourcewithref{SOURCE-DERIVED: verified private instructor-copy locators; no textbook exercise is reproduced.}{S1-53}
}

\LiveAndFullFrame{S1-34}{The analyst owns the specification and the check}{
\begin{block}{Today}
A strong target study synthesizes business configuration, operating evidence and valuation-relevant uncertainty. AI can accelerate the path from materials to a draft; it cannot decide which sources govern or whether an interpretation is justified.
\end{block}
\begin{center}
\Large\textcolor{CourseNavy}{Evidence $\rightarrow$ analysis $\rightarrow$ challenge $\rightarrow$ draft $\rightarrow$ reconciliation}
\end{center}
\begin{block}{Before next session}
Keep your group AI Work Record and final prompt-chain locator.\\ 
In Session 2, the question becomes whether a historical Excel model can be trusted after inspection and independent checks.
\end{block}
}

%\ReferenceOnlyFrame{S1-35}{Reference: master-deck release policy}{
%\begin{itemize}
%\item This full/master deck is the one substantive Session-1 source. The run of show identifies the normally projected slides; the live PDF is an automatically filtered compatibility export.
%\item The complete master deck and TargetCo profile cannot be pre-class releases while the AI-unscaffolded baseline remains part of the class. The master contains prompting method, model answer and reference workflows; the profile can be pre-solved with AI.
%\item A pre-class export, if useful, contains only non-contaminating course/logistics material from this same source. Release the complete master deck after class.
%\end{itemize}
%\sourcewithback{INSTRUCTOR PRODUCTION POLICY; conflict recorded in \texttt{SESSION1\_PRODUCTION\_MANIFEST.md}.}{S1-01}
%}

\ReferenceOnlyFrame{S1-36}{Appendix: target-study profile and later valuation work}{
\begin{columns}[T,onlytextwidth]
\column{0.48\textwidth}
\begin{block}{Business profile: why each dimension matters}
\begin{itemize}
\item Sector/subsector frames drivers, cyclicality and a starting universe.
\item Products/services distinguish commodity-like supply from specified or service-intensive offers.
\item Customers/end markets distinguish purchaser from demand exposure.
\item Channels/geography can alter commercial model, cost base and comparability.
\end{itemize}
\end{block}
\column{0.48\textwidth}
\begin{block}{Financial profile: why selected dimensions matter}
\begin{itemize}
\item Size can correlate with scale and market comparability.
\item Profitability describes conversion of sales into profit; it is not a forecast.
\item Growth needs a driver-level reading, not only a percentage.
\item Returns/credit require later data and calculations.
\end{itemize}
\end{block}
\end{columns}
\sourcewithback{Rosenbaum \& Pearl (2013), Ch. 1, Step I and Exhibit 1.3; original paraphrase.}{S1-03}
}

\ReferenceOnlyFrame{S1-37}{Appendix: what the two-page TargetCo profile can support}{
\scriptsize
\begin{tabularx}{\textwidth}{@{}p{39mm}X X@{}}
\toprule
\textbf{Evidence supplied} & \textbf{A cautious analyst may say} & \textbf{A cautious analyst must stop short of saying}\\
\midrule
End-market, geography, channel and concentration percentages & Reported mix and configuration merit follow-up. & Concentration is high/low, commercially good/bad or contractually secure.\\
Revenue bridge, utilization and margin history & 2023A price/mix contribution offsets falling volume; 2025A has a larger volume contribution; utilization and margin recover from 2024A. & A quantified causal bridge, facility margin or sustainable run-rate.\\
M-01 and M-02 management statements & What management says about recovery levers and cost classification. & That recovery will occur or an adjustment is justified.\\
\bottomrule
\end{tabularx}
\begin{block}{Classroom purpose}
The analyst must synthesize, calculate, challenge and prioritize. The short source pack is deliberately incomplete, not a hidden answer.
\end{block}
\sourcewithback{Back.}{S1-10}
%\sourcewithback{TARGETCO APPLICATION; INSTRUCTOR SYNTHESIS.}{S1-10}
}

\ReferenceOnlyFrame{S1-38}{Appendix: four-dimensional rubric calibration}{
\scriptsize
\begin{tabularx}{\textwidth}{@{}p{29mm}p{30mm}p{30mm}X@{}}
\toprule
\textbf{Dimension} & \textbf{0} & \textbf{1} & \textbf{2}\\
\midrule
Business-model synthesis & Copies labels or omits configuration. & Names facts but weakly connects them. & Connects specification/qualification, customer/end market, route to market and footprint without inventing economics.\\
Historical economics & Wrong contribution story, arithmetic or causality. & Lists figures with a weak bridge. & Reconciles growth, distinguishes contributions, treats utilization/margin cautiously and keeps numerical fidelity.\\
Evidence and uncertainty discipline & Imports facts or treats management assertion as fact. & Notes uncertainty but blurs categories. & Separates observed data, management statements and qualified inference; turns unsupported claims into questions.\\
Prioritization / VP usefulness & Not three questions or generic list. & Three questions but weak consequence or unclear note. & Exactly three ranked questions; each names a later judgement; concise VP-ready note.\\
\bottomrule
\end{tabularx}
\sourcewithback{Back.}{S1-09}
%\sourcewithback{INSTRUCTOR SYNTHESIS: score actual output, not presumed model quality.}{S1-09}
}

\ReferenceOnlyFrame{S1-39}{Appendix: model, application and tool in the TargetCo workflow}{
\small
\begin{block}{Concrete pathway}
TargetCo PDF $\rightarrow$ application parses the file and manages the chat/context $\rightarrow$ model generates the response $\rightarrow$ a calculator/code/file tool may retrieve, recompute or check a result $\rightarrow$ analyst reviews it.
\end{block}
\begin{tabularx}{\textwidth}{@{}p{30mm}X@{}}
\toprule
\textbf{Term} & \textbf{Operational consequence}\\
\midrule
Model & Produces conditional next-token predictions. Record any displayed model/version; do not infer hidden configuration.\\
Application & May add file conversion, system instructions, context management, routing, memory or tool access. An application comparison is not a pure model experiment.\\
Tool & Retrieves, calculates or changes an external object. Tool output creates a source and checking obligation.\\
\bottomrule
\end{tabularx}
\sourcewithback{Back.}{S1-16}
%\sourcewithback{Official OpenAI and Anthropic product/prompt guidance.}{S1-16}
}

\ReferenceOnlyFrame{S1-40}{Appendix: long-context failure modes}{
\footnotesize
\begin{columns}[T,onlytextwidth]
\column{0.48\textwidth}
\key{Failure mode}
\begin{itemize}
\item management material and an authoritative source are treated as equally established;
\item source material contains instructions that compete with the task;
\item material evidence is lost among irrelevant pages;
\item a citation label appears without a reconciled claim.
\end{itemize}
\column{0.48\textwidth}
\key{Design response}
\begin{itemize}
\item label source, date and authority;
\item delimit documents and permitted use;
\item make the task instruction easy to locate;
\item extract fact/statement/calculation with location before synthesis.
\end{itemize}
\end{columns}
\begin{block}{TargetCo lesson}
The profile is authoritative for the exercise while M-01/M-02 within it remain management statements. Authority and truth status are different analytical questions.
\end{block}
\sourcewithback{Back.}{S1-15}
%\sourcewithback{Anthropic Prompting Best Practices (long-context section).}{S1-15}
}

\ReferenceOnlyFrame{S1-41}{Appendix: evaluation comes before prompt cleverness}{
\begin{block}{Evaluation loop}
\[
\text{task-specific test cases} \rightarrow \text{initial prompt} \rightarrow \text{run} \rightarrow \text{score} \rightarrow \text{diagnose} \rightarrow \text{revise}.
\]
\end{block}
\scriptsize
\begin{tabularx}{\textwidth}{@{}p{43mm}X X@{}}
\toprule
\textbf{Observed failure} & \textbf{Precise diagnosis} & \textbf{Response}\\
\midrule
Calls 2023A demand/volume-led & The output ignores disclosed $-4.0$ ppt volume and +6.1 ppt price/mix contributions. & Require source-grounded extraction and contribution recomputation before prose.\\
Treats M-02 as fact & Management-statement classification/source-fidelity check is missing. & Request classification, locator and analyst verification before drafting.\\
Gives six generic questions & Prioritization criterion is not operational. & Require exactly three ranked questions and one named valuation consequence each.\\
Correct but too long & Output constraint fails. & State one-page constraint and use a non-TargetCo output-shape example if format remains unstable.\\
\bottomrule
\end{tabularx}
\sourcewithback{Back.}{S1-18}
%\sourcewithback{Anthropic Define Success Criteria and Build Evaluations.}{S1-18}
}

\ReferenceOnlyFrame{S1-42}{Appendix: few-shot examples and the baseline boundary}{
\begin{block}{Strong non-TargetCo output-shape example}
\textbf{Observed fact:} a product requires customer qualification before serial production.\par
\textbf{Restrained interpretation:} qualification may make switching and implementation relevant to diligence.\par
\textbf{Open question:} what evidence shows the duration, economics and renewal pattern of qualified relationships?
\end{block}
\begin{columns}[T,onlytextwidth]
\column{0.49\textwidth}
\key{Appropriate}
Use this type of example after baseline when a group needs a model of fact $\rightarrow$ qualified interpretation $\rightarrow$ question.
\column{0.49\textwidth}
\key{Inappropriate before baseline}
A filled TargetCo example leaks the analytical answer, anchors wording and turns the task into transcription.
\end{columns}
\sourcewithback{Back.}{S1-22}
%\sourcewithback{SOURCE-DERIVED: Anthropic Prompting Best Practices, use examples effectively; INSTRUCTOR SYNTHESIS.}{S1-22}
}

\ReferenceOnlyFrame{S1-43}{Appendix: the same semantic prompt in headings and XML}{
\scriptsize
\begin{columns}[T,onlytextwidth]
\column{0.47\textwidth}
\key{Neutral headings}\par\smallskip
\ttfamily
Task\newline
Prepare a one-page TargetCo VP note.\newline
Sources\newline
Supplied TargetCo profile only.\newline
Output\newline
Four named sections.\newline
Criteria\newline
Four observable success criteria.\newline
Checks\newline
Recompute bridge; reconcile claims.
\column{0.50\textwidth}
\key{Anthropic XML}\par\smallskip
\ttfamily
\textless task\textgreater Prepare a one-page TargetCo VP note. \textless/task\textgreater\newline
\textless sources\textgreater Supplied TargetCo profile only. \textless/sources\textgreater\newline
\textless output\textgreater Four named sections. \textless/output\textgreater\newline
\textless criteria\textgreater Four observable success criteria. \textless/criteria\textgreater\newline
\textless checks\textgreater Recompute bridge; reconcile claims. \textless/checks\textgreater
\end{columns}
\smallskip
The tags separate task, source boundary, requested output, criteria and checks. They do not create a different reasoning method or independently verify a result.
\sourcewithback{Back.}{S1-23}
%\sourcewithback{SOURCE-DERIVED: Anthropic Prompting Best Practices, structure prompts with XML tags; INSTRUCTOR SYNTHESIS.}{S1-23}
}

\ReferenceOnlyFrame{S1-44}{Appendix: roles set audience and standards}{
\footnotesize
\begin{columns}[T,onlytextwidth]
\column{0.48\textwidth}
\begin{block}{Prompt A}
You are a senior investment banker. Explain TargetCo's prospects and give a valuation view.
\end{block}
\risk{Problem:} the role conflicts with the assignment, invites external knowledge and encourages an unsupported valuation claim.
\column{0.48\textwidth}
\begin{block}{Prompt B}
Write for a VP considering a sell-side mandate. Use only the supplied profile. State the operating picture, label management statements and leave unsupported causality as questions. Do not value the company.
\end{block}
\good{Benefit:} audience, vocabulary, source boundary and standards are explicit.
\end{columns}
\begin{alertblock}{Rule}
A role guides tone, context and perspective. It cannot add facts, make an inference true or transfer accountability.
\end{alertblock}
\sourcewithback{Back.}{S1-22}
%\sourcewithback{SOURCE-DERIVED: Anthropic Prompting Best Practices, give Claude a role; INSTRUCTOR SYNTHESIS.}{S1-22}
}

\ReferenceOnlyFrame{S1-45}{Appendix: prompt-chain skeleton}{
\scriptsize
\begin{center}
\begin{tabular}{p{20mm}p{31mm}p{44mm}p{24mm}}
\toprule
\textbf{Stage} & \textbf{Output} & \textbf{Human question} & \textbf{Failure localized}\\
\midrule
Extract & Evidence table; calculations & Right facts, contribution points and source locations? & Retrieval / arithmetic\\
Challenge & Criteria-based corrections & What is unsupported, omitted or generic? & Interpretation / specification\\
Draft & One-page VP note & Is it concise, qualified and decision-useful? & Communication / prioritization\\
Verify & Reconciliation log & Does each material point return to a source or calculation? & Source fidelity\\
\bottomrule
\end{tabular}
\end{center}
\begin{block}{State-management cost}
Each stage needs a stable source pack, versioned prompt and handoff. The live redo compresses the four logical stages into at most three user messages.
\end{block}
\sourcewithback{Back.}{S1-19}
%\sourcewithback{SOURCE-DERIVED: Anthropic Prompt Engineering Overview and Prompting Best Practices; INSTRUCTOR SYNTHESIS.}{S1-19}
}

\ReferenceOnlyFrame{S1-46}{Appendix: self-check versus independent check}{
\small
\begin{tabularx}{\textwidth}{@{}p{31mm}p{48mm}X@{}}
\toprule
\textbf{Check request} & \textbf{What it can reasonably do} & \textbf{What is still required}\\
\midrule
``Review this draft against the rubric.'' & Detect missing sections, generic questions or an obvious internal contradiction. & Reconcile each material statement to the supplied profile.\\
``Check the calculations.'' & Recalculate a percentage or identify arithmetic inconsistency. & Independently recompute in calculator/spreadsheet and identify inputs/source.\\
``Are you sure M-02 is non-recurring?'' & Surface qualification language or admit missing support. & Obtain evidence on cost nature/duration; analyst decides later normalization treatment.\\
\bottomrule
\end{tabularx}
\begin{alertblock}{TargetCo implication}
M-02 spans two years. That supports a question about recurrence and duration; it neither proves recurrence nor non-recurrence, and it establishes no adjustment amount.
\end{alertblock}
\sourcewithback{Back.}{S1-24}
%\sourcewithback{INSTRUCTOR SYNTHESIS; TARGETCO APPLICATION.}{S1-24}
}

\ReferenceOnlyFrame{S1-47}{Appendix: fuller TargetCo evidence matrix}{
\scriptsize
\begin{tabularx}{\textwidth}{@{}p{25mm}p{24mm}p{32mm}X@{}}
\toprule
\textbf{Analytical point} & \textbf{Source location} & \textbf{Evidence/calculation} & \textbf{Support, limit and next question}\\
\midrule
Specification-led configuration & Company extract; product/facility notes & Qualification before serial production; containers, closures and dispensing components; 71\% direct/key account mix & Supports specialty manufacturing/service configuration; does not show pricing power or channel margin.\\
2023A growth character & Historical table; 2023A commentary & Volume contribution $-4.0$ ppt; price/mix +6.1 ppt; revenue +2.1\% & Supports contrast with 2025A; no customer-level price/mix bridge.\\
2025A utilization/margin tension & 2022A--2025A table & Utilization 82\% to 84\%; margin 14.8\% to 14.3\% & Supports challenge to single-driver causality; open cost/mix/facility bridge.\\
Customer concentration & Commercial snapshot & Largest 9.4\%; top ten 41.8\% & Supports request for contract and retention context; no automatic risk conclusion.\\
M-01/M-02 & Management materials & Asserted recovery levers; two years of stated non-recurring costs & Supports diligence questions; does not establish forecast or adjustment.\\
\bottomrule
\end{tabularx}
\sourcewithback{Back.}{S1-26}
%\sourcewithback{INSTRUCTOR SYNTHESIS; TARGETCO APPLICATION.}{S1-26}
}

\ReferenceOnlyFrame{S1-48}{Appendix: why the three questions are ranked}{
\begin{enumerate}
\item \textbf{Margin bridge and supported run-rate.} The utilization/margin tension leaves an unobserved bridge. The answer could alter margin forecasts and valuation.
\item \textbf{Nature and duration of commissioning/quality costs.} Identifiable, finite, incremental and non-run-rate evidence affects whether any normalization is supportable. Two-year duration creates a challenge; it does not establish a treatment.
\item \textbf{Breadth and durability of volume recovery.} Evidence across customer, channel, end market and facility can change forecast confidence, risk and peer selection.
\end{enumerate}
\begin{block}{Acceptable alternative}
A question about qualification timing or direct-account/distributor economics can be equally strong if it states a potential valuation consequence and is prioritized against these questions.
\end{block}
\sourcewithback{Back.}{S1-29}
%\sourcewithback{TARGETCO APPLICATION; INSTRUCTOR SYNTHESIS.}{S1-29}
}

\ReferenceOnlyFrame{S1-49}{Appendix: instructor prompt-chain skeleton}{
\scriptsize
\begin{block}{Message 1: evidence / analysis}
Use only the supplied TargetCo material. Extract business-model evidence; recompute revenue growth; distinguish disclosed volume and price/mix contribution points; track utilization and margin; identify tensions; return source location, support and limit. No valuation or final prose.
\end{block}
\begin{block}{Message 2: challenge + draft}
Check the analysis against the four criteria. Correct unsupported inference, omitted dimensions, numerical inconsistency and generic questions. Then draft the one-page note under Business model; Historical trajectory; Key tensions; Three ranked questions + why.
\end{block}
\begin{block}{Message 3: verify + revise}
Reconcile each material number and statement to source or calculation. Revise the note only for verified corrections. Flag remaining uncertainty rather than supplying a fact.
\end{block}
\sourcewithback{Back.}{S1-30}
%\sourcewithback{INSTRUCTOR SYNTHESIS: full near-final version follows; full text in \texttt{instructor/SESSION1\_INSTRUCTOR\_AI\_WORKFLOW.md}.}{S1-30}
}

\ReferenceOnlyFrame{S1-50}{Appendix: a strong near-final prompt sequence for TargetCo}{
\scriptsize
\textbf{Message 1 -- evidence / analysis.} 
The supplied TargetCo profile is authoritative. Use no external facts, tools or valuation. Extract product/qualification, markets, channel and footprint; revenue and $g_t=R_t/R_{t-1}-1$; disclosed volume and price/mix contribution points (check $g_t=v_t+p_t$); utilization and reported margin; and M-01/M-02 as management statements. For each, give profile location, classification, calculation/support and limit. Flag 2025A's 84\% versus 82\% utilization but 14.3\% versus 14.8\% margin tension. Do not draft.

\smallskip
\textbf{Message 2 -- challenge + draft.} 
Assess the extraction against four criteria: coherent business-model synthesis; correct contribution arithmetic and qualified causality; evidence/uncertainty discipline; concise VP usefulness with exactly three ranked questions naming a later judgement. Correct only what the profile supports. Draft one page: Business model / positioning; Historical trajectory; Key tensions; Three ranked questions + why. Preserve M-01/M-02 as management statements. No valuation.

\smallskip
\textbf{Message 3 -- verify + revise.} 
Reconcile every material number and statement to a profile location or calculation. Recheck 2023A ($-4.0$ ppt + 6.1 ppt = 2.1\%) and 2025A (4.0 ppt + 1.9 ppt = 5.9\%). Confirm that no management statement became fact, causal language remains qualified and exactly three questions are ranked. Give a short reconciliation list, revise only where required, and flag unresolved uncertainty.
\sourcewithback{Back.}{S1-30}
%\sourcewithback{INSTRUCTOR SYNTHESIS: a strong near-final sequence, not an objectively perfect prompt.}{S1-30}
}

\ReferenceOnlyFrame{S1-51}{Appendix: common weak moves are diagnostic}{
\footnotesize
\begin{tabularx}{\textwidth}{@{}p{42mm}X@{}}
\toprule
\textbf{Weak move} & \textbf{Why it fails and how to repair it}\\
\midrule
Rewrites profile headings & It has not synthesized an operating configuration. Connect product, qualification, channel and geography cautiously.\\
Lists annual data & It misses changing economics. Contrast 2023A price/mix contribution offsetting falling volume with 2025A's larger volume contribution.\\
Says utilization ``caused'' margin recovery & The pack has no quantified causal bridge. Use qualified language and ask for the bridge.\\
Calls M-01/M-02 facts & It converts management language into established evidence. Label the statement and specify missing support.\\
Calls 9.4\% / 41.8\% high or low & Benchmark and contractual context are absent. State the figures and ask for context.\\
Asks generic diligence questions & A VP cannot see valuation relevance. Name the later judgement that might change.\\
\bottomrule
\end{tabularx}
\sourcewithback{Back.}{S1-31}
%\sourcewithback{INSTRUCTOR SYNTHESIS; TARGETCO APPLICATION.}{S1-31}
}

\ReferenceOnlyFrame{S1-52}{Appendix: Session 1 and the AI Fluency 4Ds}{
\scriptsize
\begin{tabularx}{\textwidth}{@{}p{22mm}p{47mm}X@{}}
\toprule
\textbf{4D} & \textbf{Session 1 action} & \textbf{Evidence retained}\\
\midrule
Delegation & Delegate extraction, cross-check suggestions and drafting; retain authority, valuation relevance and approval. & Work Record: delegated versus retained.\\
Description & State task, authoritative source, output, criteria and checks. & P0 and P1--P4 versions.\\
Discernment & Score four dimensions; identify unsupported stories and omissions. & Rubric, challenge, baseline/redo.\\
Diligence & Reconcile sources; disclose uncertainty; do not treat management language as verified. & Matrix, numerical check, remaining uncertainty.\\
\bottomrule
\end{tabularx}
\begin{block}{Mode}
This is augmentation: AI extends analysis under human direction. Later agency needs scope, approvals, state and stopping rules.
\end{block}
\sourcewithback{Back.}{S1-17}
%\sourcewithback{SOURCE-DERIVED: Anthropic AI Fluency: Framework \& Foundations; INSTRUCTOR SYNTHESIS.}{S1-17}
}

\ReferenceOnlyFrame{S1-53}{Appendix: finance locators and boundaries}{
\footnotesize
\begin{block}{Rosenbaum \& Pearl (2013), \emph{Investment Banking}}
Ch. 1, Step I ``Study the Target''; Exhibit 1.3 Business and Financial Profile Framework; ``Financial Profile''; and ``Calculation of Key Financial Statistics and Ratios.'' Used for target-study sequence, profile dimensions and Session-2 navigation.
\end{block}
\begin{block}{Rosenbaum \& Pearl Workbook (2013) and Koller et al. (2025)}
Workbook Ch. 1 Q14--21 (business/financial profile), Q7--9 (growth/profitability/credit calculations), Q3 (one-time/non-recurring adjustments); Koller Chs. 11 and 13 for optional model-organisation/hygiene context.
\end{block}
\begin{alertblock}{Rights boundary}
Slides paraphrase methods and cite locators. They do not reproduce protected wording, tables, exhibits, cases, exercises, answers, layouts or values.
\end{alertblock}
\sourcewithback{Back.}{S1-33}
%\sourcewithback{SOURCE-DERIVED: private instructor-held copies consulted for locator alignment.}{S1-33}
}

\ReferenceOnlyFrame{S1-54}{Appendix: sources, claims and application}{
\begin{tabularx}{\textwidth}{@{}p{37mm}X@{}}
\toprule
\textbf{Label} & \textbf{Meaning in this deck}\\
\midrule
SOURCE-DERIVED & A principle is paraphrased from an identified official page or a textbook section/exhibit locator.\\
INSTRUCTOR SYNTHESIS & A teaching framework, rubric, prompt chain, explanation or wording designed for this course. It is not presented as a source quotation.\\
TARGETCO APPLICATION & The source-derived or instructor-synthesis method is applied only to disclosed canonical TargetCo facts and labelled management statements.\\
\bottomrule
\end{tabularx}
\begin{block}{Source map}
\texttt{SLIDE\_SOURCE\_MAP.md} records exact locators and adopted principles by internal slide anchor. \texttt{SESSION1\_SOURCE\_NOTES.md} records source access and rights boundaries. The map enables review; it does not replace source reading.
\end{block}
\sourcewithback{Back.}{S1-01}
%\sourcewithback{INSTRUCTOR SYNTHESIS: source-governance convention for Session 1.}{S1-01}
}
